\section*{第 \ref{ch:b1:layered-earth} 章——分层地球}

\begin{solution}{exo:b1:layered-earth:1}
体积分数等于半径比的立方：地核
$(3480/6371)^3 = 0.1630$，即 \qty{16.3}{\percent}；\omterm{def:b1:layered-earth:layers}{内核}
$(1221.5/6371)^3 = 0.00705$, \qty{0.70}{\percent}.
\end{solution}

\begin{solution}{exo:b1:layered-earth:2}
上方：$\Phi = 13.7166^2 - \tfrac43\times7.26466^2 = 188.145 - 70.367 =
\qty{117.78}{km^2.s^{-2}}$，$K = \rho\Phi = \qty{5566.45}{kg.m^{-3}}\times
\qty{117.78e6}{m^2.s^{-2}} = \qty{655.6}{GPa}$。下方：$\Phi = 8.06482^2 =
\qty{65.04}{km^2.s^{-2}}$，$K = 9903.49\times\num{65.04e6} = \qty{644.1}{GPa}$。
体积模量几乎不变（\qty{-1.7}{\percent}）。变化的是
剪切模量 $\mu = \rho v_S^2 = \qty{293.8}{GPa}$ 在上方，下方则为 $0$（
液体），密度则上升 \qty{78}{\percent}，这就解释了
$v_P$ 下降的原因，尽管 $K$ 几乎相等。
\end{solution}

\begin{solution}{exo:b1:layered-earth:3}
$\sin i_c = 6.2/8.1 = 0.7654$，$i_c = \ang{49.95}$。临界距离
$2\times30\times\tan\ang{49.95} = 60\times1.1906 = \qty{71.4}{km}$。\omterm{def:b1:layered-earth:head-wave}{交叉
距离} $60\sqrt{14.3/1.9} = 60\times2.743 = \qty{164.6}{km}$。
\end{solution}

\begin{solution}{exo:b1:layered-earth:4}
$g = Gm/r^2 = \num{6.674e-11}\times\num{1.94e24}/(\num{3.48e6})^2 =
\qty{10.69}{m.s^{-2}}$，比地表高 \qty{9}{\percent}。向下
至地核时，半径平方降为地表值的 $(3480/6371)^2 = 0.298$，
而吸引质量仅降为总质量的 $1.94/5.972 = 0.325$，
因为地核约占质量的三分之一：$0.325/0.298 = 1.09$。
\end{solution}

\begin{solution}{exo:b1:layered-earth:5}
靠近震源处 $v_1 = 60/10.0 = \qty{6.0}{km.s^{-1}}$。远处
$v_2 = (300 - 200)/(45.8 - 33.3) = 100/12.5 = \qty{8.0}{km.s^{-1}}$，截距时间
为 $t_i = 33.3 - 200/8.0 = \qty{8.3}{s}$。然后 $\sin i_c =
0.75$，$\cos i_c = 0.6614$，则 $h = t_i v_1/(2\cos i_c) = 8.3\times6.0/1.3229
= \qty{37.6}{km}$。验算：$x_c = 2\times37.6\times\sqrt{14/2} = \qty{199}{km}$，
两种波都在 $199/6.0 = 199/8.0 + 8.3 = \qty{33.2}{s}$ 时刻到达。
\end{solution}

\begin{solution}{exo:b1:layered-earth:6}
对于均匀的 $\rho$，$I = \tfrac{8\pi}{3}\rho\int_0^R r^4\,\mathrm dr =
\tfrac{8\pi}{15}\rho R^5$，且 $M = \tfrac43\pi\rho R^3$，故
$I/MR^2 = (8/15)/(4/3) = 0.4$。两层（$r_c = \qty{3.48e6}{m}$，
$R = \qty{6.371e6}{m}$）：$M = \tfrac43\pi\bigl(11000\,r_c^3 + 4450\,(R^3 -
r_c^3)\bigr) = \qty{5.977e24}{kg}$，比实测
质量高 \qty{0.08}{\percent}；$I = \tfrac{8\pi}{15}\bigl(11000\,r_c^5 + 4450\,(R^5 - r_c^5)\bigr)$ 给出
$I/MR^2 = 0.3457$。质量正确，但该因子偏高了 \qty{4.5}{\percent}
：实际\omterm{def:b1:layered-earth:layers}{地幔}并非均匀，其密度随深度增大，
这使质量向内迁移。
\end{solution}

\begin{solution}{exo:b1:layered-earth:7}
$\Phi = 12.16^2 - \tfrac43\times6.673^2 = 147.87 - 59.37 =
\qty{88.49}{km^2.s^{-2}}$。以 $z$ 表示深度，
$\mathrm d\rho/\mathrm dz = \rho g/\Phi = 4860\times9.9/\num{88.49e6} =
\qty{5.44e-4}{kg.m^{-4}}$，即每千米 \qty{0.544}{kg.m^{-3}}，或
\qty{0.54}{g.cm^{-3}}，每 \qty{1000}{km}。PREM 增加 $5.357 - 4.844 =
0.513$，在这段 \qty{1000}{km} 上少了 \qty{6}{\percent}：
在顶部算得的梯度被用于整个区间，而 $\Phi$ 随深度增大，
梯度随之减小。
\end{solution}

\begin{solution}{exo:b1:layered-earth:8}
$P \approx \bar\rho g z = 3570\times9.9\times\num{670e3} = \qty{23.7}{GPa}$
（PREM：\qty{23.9}{GPa}）。这一压强约为 \qty{24}{GPa}，
在此压强下林伍德石分解为布里奇曼石与铁方镁石，即
\qty{660}{km} 不连续面。
\end{solution}

\begin{solution}{exo:b1:layered-earth:9}
(a)~液态铁：$10.90\times1.05 = 11.4$ 至 $10.90\times1.10 =
\qty{12.0}{g.cm^{-3}}$。(b)~$12.764/12.166 = 1.049$：PREM 中的跃变
为 \qty{4.9}{\percent}，略高于
近期研究的 \qty{4.5}{\percent}。(c)~仅凝固一项，只会使固体比液体
稍微致密一些；跃变更大，说明固体所含的轻元素比液体少：
轻元素更\omterm{def:b1:geological-maps:attitude}{倾向}于留在液体中，因此在
\omterm{def:b1:layered-earth:discontinuity}{内核边界}凝固时，大部分轻元素被留了下来。
\end{solution}

\begin{solution}{exo:b1:layered-earth:10}
设 $x_* = 2h\tan i_c$。则 $\sqrt{x_*^2 + 4h^2} = 2h/\cos i_c$，且
$t_r(x_*) = 2h/(v_1\cos i_c)$。\omterm{def:b1:layered-earth:head-wave}{首波}给出
$t(x_*) = 2h\tan i_c/v_2 + 2h\cos i_c/v_1 = \frac{2h}{v_1}\bigl(\sin^2 i_c/\cos
i_c + \cos i_c\bigr) = 2h/(v_1\cos i_c)$，利用 $\sin i_c = v_1/v_2$，两者相同。
双曲线的斜率为 $t_r'(x) = x/(v_1\sqrt{x^2 + 4h^2})$；在 $x_*$ 处为
$\sin i_c/v_1 = 1/v_2$，即\omterm{def:b1:layered-earth:head-wave}{首波}直线的斜率。由于
$t_r'' > 0$，双曲线是凸的，位于其每条切线之上：对于
$x > x_*$ 时，$t_r > t_{\text{首波}}$。在临界距离处，反射
恰以 $i_c$ 射到界面，此处即\omterm{def:b1:layered-earth:head-wave}{首波}的诞生点；超过此距离后，
\omterm{def:b1:layered-earth:head-wave}{首波}因在快速介质中传播而得益，先行到达。
\end{solution}

\begin{solution}{exo:b1:layered-earth:11}
$m(r) = \tfrac43\pi\rho r^3$，故 $g = Gm/r^2 = \tfrac43\pi G\rho r$。于是
$P(r) = \int_r^R \rho\,\tfrac43\pi G\rho s\,\mathrm ds = \tfrac23\pi G\rho^2(R^2
- r^2)$。地心处 $P = \tfrac23\pi\times\num{6.674e-11}\times5513^2\times
(\num{6.371e6})^2 = \qty{172}{GPa}$，约为 PREM 值
\qty{363.85}{GPa} 的 \qty{47}{\percent}。把质量向中心集中，
大致使中心压强翻了一倍。
\end{solution}

\begin{solution}{exo:b1:layered-earth:12}
设 $z$ 为深度，$\mathrm d\rho/\mathrm dz = \rho g/\Phi$；若 $g/\Phi$ 为
常数 $a$，则 $\mathrm d(\ln\rho)/\mathrm dz = a$，且 $\rho = \rho_0
e^{a(z - z_0)}$。此处 $a(z - z_0) = 7.9\times\num{2258.5e3}/\num{87e6} =
0.2051$，则 $\rho = 9.90\times e^{0.2051} = 9.90\times1.2276 =
\qty{12.15}{g.cm^{-3}}$，比 PREM 的 $12.17$ 低 \qty{0.2}{\percent}：\omterm{def:b1:layered-earth:layers}{外核}内
密度的全部增加都是自压缩的结果。
\end{solution}

\begin{solution}{pb:b1:layered-earth:1}
\textbf{1.} $V = \tfrac43\pi(\num{6.371e6})^3 = \qty{1.0832e21}{m^3}$.

\textbf{2.} $\tfrac43\pi(r_2^3 - r_1^3)$：\omterm{def:b1:layered-earth:layers}{地壳}、上\omterm{def:b1:layered-earth:layers}{地幔}和\omterm{def:b1:layered-earth:transition-zone}{地幔过渡带}
\qty{3.0706e20}{m^3}；下\omterm{def:b1:layered-earth:layers}{地幔} \qty{5.9961e20}{m^3}；\omterm{def:b1:layered-earth:layers}{外核}
\qty{1.6890e20}{m^3}；\omterm{def:b1:layered-earth:layers}{内核} \qty{7.634e18}{m^3}。

\textbf{3.} \qty{28.35}{\percent}, \qty{55.36}{\percent}, \qty{15.59}{\percent}
和 \qty{0.70}{\percent}：下\omterm{def:b1:layered-earth:layers}{地幔}是最大的壳层，
超过地球体积的一半。

\textbf{4.} $15.59 + 0.70 = \qty{16.3}{\percent}$.

\textbf{5.} \omterm{def:b1:layered-earth:layers}{地壳}、上\omterm{def:b1:layered-earth:layers}{地幔}和\omterm{def:b1:layered-earth:transition-zone}{地幔过渡带} $\num{3.0706e20}\times3570 =
\qty{1.0962e24}{kg}$；下\omterm{def:b1:layered-earth:layers}{地幔} \qty{2.9381e24}{kg}；\omterm{def:b1:layered-earth:layers}{外核}
\qty{1.8410e24}{kg}；\omterm{def:b1:layered-earth:layers}{内核} \qty{9.841e22}{kg}。

\textbf{6.} 质量为 \qty{5.9737e24}{kg}，比实测值
\qty{5.972e24}{kg} 高 \qty{0.03}{\percent}。

\textbf{7.} $\num{5.9737e24}/\num{1.0832e21} = \qty{5515}{kg.m^{-3}}$.

\textbf{8.} $\num{9.841e22}/\num{5.9737e24} = \qty{1.65}{\percent}$.

\textbf{9.} $(1.0962 + 2.9381)/5.9737 = \qty{67.5}{\percent}$.

\textbf{10.} 地核体积占 \qty{16.3}{\percent}，质量占
\qty{32.5}{\percent}，是体积占比的两倍：其平均密度
\qty{10986}{kg.m^{-3}}，约为\omterm{def:b1:layered-earth:layers}{地幔}的 $2.5$ 倍。

\textbf{11.} 密度均匀的壳层给出
$I = \tfrac{8\pi}{3}\rho\int_{r_1}^{r_2} r^4\,\mathrm dr = \tfrac{8\pi}{15}\rho
(r_2^5 - r_1^5)$.

\textbf{12.} \omterm{def:b1:layered-earth:layers}{地壳}、上\omterm{def:b1:layered-earth:layers}{地幔}和\omterm{def:b1:layered-earth:transition-zone}{地幔过渡带} \qty{2.676e37}{kg.m^2}；下\omterm{def:b1:layered-earth:layers}{地幔}
\qty{4.525e37}{kg.m^2}；\omterm{def:b1:layered-earth:layers}{外核} \qty{9.27e36}{kg.m^2}；\omterm{def:b1:layered-earth:layers}{内核}
\qty{5.87e34}{kg.m^2}。

\textbf{13.} 下\omterm{def:b1:layered-earth:layers}{地幔}：它占一半质量，且位于较大半径处。\omterm{def:b1:layered-earth:layers}{内
核}占质量的 \qty{1.6}{\percent}，并全部靠近转轴：
其占比为 \qty{0.07}{\percent}。

\textbf{14.} $I = \qty{8.1345e37}{kg.m^2}$ 及 $I/MR^2 = \num{8.1345e37}/
(\num{5.9737e24}\times(\num{6.371e6})^2) = 0.3355$.

\textbf{15.} 比 $0.3308$ 高出 \qty{1.4}{\percent}，远低于 $0.4$：
模型对质量向中心集中程度的描述基本正确。

\textbf{16.} 是。每个壳层内部，真实密度在底部
高于顶部；质量相同的均匀壳层会把过多质量
集中在外半径附近，该处以 $r^2$ 计入，因而高估 $I$。

\textbf{17.} $g = GM/R^2 = \num{6.674e-11}\times\num{5.9737e24}/
(\num{6.371e6})^2 = \qty{9.822}{m.s^{-2}}$.

\textbf{18.} 地核质量为 $\num{1.8410e24} + \num{9.841e22} =
\qty{1.9394e24}{kg}$，且 $g = G m/r_c^2 = \qty{10.69}{m.s^{-2}}$。

\textbf{19.} $(\num{1.0962e24} + \num{2.9381e24})/(\num{3.0706e20} +
\num{5.9961e20}) = \num{4.0343e24}/\num{9.0667e20} = \qty{4450}{kg.m^{-3}}$.

\textbf{20.} $P \approx 4450\times10\times\num{2.891e6} = \qty{128.6}{GPa}$,
比 PREM 的 \qty{135.75}{GPa} 低 \qty{5}{\percent}。

\textbf{21.} $\rho_c = \num{1.9394e24}/\num{1.7653e20} = \qty{10986}{kg.m^{-3}}$,
且 $\Delta P = \tfrac23\pi\times\num{6.674e-11}\times10986^2\times
(\num{3.48e6})^2 = \qty{204.3}{GPa}$.

\textbf{22.} $128.6 + 204.3 = \qty{332.9}{GPa}$，比
PREM 的 \qty{363.85}{GPa} 低 \qty{8.5}{\percent}。两种估算都忽略了各区域内部密度随
深度增大的影响，这使更多质量位于深处，并使重力在向下过程中保持较高，
从而提高 $\rho g$，即压强梯度。

\textbf{23.} 总质量（第 6 题）、\omterm{def:b1:comparative-planetology:moi}{转动惯量因子}
（第 15 题）和地表重力（第 17 题，对照实测平均值
\qty{9.82}{m.s^{-2}}）均可用地球的测量结果检验；地核处的重力
和压强只能与 PREM 比较，而 PREM 本身
也只是一个模型。

\textbf{24.} 地核占地球质量的
\textbf{\qty[text-series-to-math=true]{32.5}{\percent}}。
\end{solution}
